% TV1 (Tạ Bỉnh Phong - 24127100) viết file này; không sửa định dạng chung.
% Tương ứng các mục 1--2 và phần dữ liệu/baseline trong project plan.
\chapter{Giới thiệu và các hướng tiếp cận trước}
\label{ch:introduction}
\owner{TV1 -- \MemberOne}{TV2 và TV5}

\section{Bối cảnh và phát biểu bài toán}
Bài toán hồi quy và xấp xỉ hàm nhiều biến là một trong những nhiệm vụ nền tảng của học máy có giám sát. Cho tập quan sát hữu hạn
\begin{equation}\label{eq:dataset}
\mathcal{D} = \left\{(\mathbf{x}_i, y_i)\right\}_{i=1}^{n}, \qquad \mathbf{x}_i \in \Omega \subset \mathbb{R}^{d}, \quad y_i \in \mathbb{R},
\end{equation}
trong đó $\mathbf{x}_i = (x_{i,1}, \dots, x_{i,d})^\top$ là véc-tơ đặc trưng đầu vào thuộc miền xác định $\Omega$ có số chiều $d$, còn $y_i = f(\mathbf{x}_i) + \varepsilon_i$ là giá trị mục tiêu sinh ra bởi hàm ẩn $f : \Omega \to \mathbb{R}$ (với $\varepsilon_i = 0$ trong các bài toán tổng hợp không nhiễu). Mục tiêu của mô hình tham số $f_\theta : \Omega \to \mathbb{R}$ (tham số hóa bởi $\theta \in \Theta$) là tìm bộ tham số tối ưu $\theta^\star$ sao cho $f_{\theta^\star}(\mathbf{x})$ xấp xỉ chính xác $f(\mathbf{x})$ trên toàn miền $\Omega$, kể cả tại các điểm chưa quan sát trong tập kiểm tra. Trong huấn luyện, hàm mất mát sai số bình phương trung bình (Mean Squared Error -- MSE) được tối thiểu hóa:
\begin{equation}\label{eq:intro-mse}
\mathcal{L}_{\mathrm{MSE}}(\theta) = \frac{1}{n_{\mathrm{train}}}\sum_{i=1}^{n_{\mathrm{train}}}\left(f_\theta(\mathbf{x}_i) - y_i\right)^2,
\end{equation}
và chất lượng tổng quát hóa trên một tập dữ liệu đánh giá kích thước $n_{\mathrm{eval}}$ được đo bằng căn bậc hai sai số bình phương trung bình ($\RMSE$):
\begin{equation}\label{eq:intro-rmse}
\RMSE = \sqrt{\frac{1}{n_{\mathrm{eval}}}\sum_{i=1}^{n_{\mathrm{eval}}}\left(f_\theta(\mathbf{x}_i) - y_i\right)^2}.
\end{equation}

Trong nhiều thập kỷ qua, Multi-Layer Perceptron (MLP) là kiến trúc mặc định cho bài toán xấp xỉ hàm phi tuyến nhờ định lý xấp xỉ vạn năng~\cite{cybenko1989approximation,hornik1989multilayer}. Tuy nhiên, khi ứng dụng vào các bài toán khoa học và toán học quy mô nhỏ cần độ chính xác số cao hoặc khả năng diễn giải công thức, Liu và cộng sự~\cite{liu2025kan} chỉ ra rằng MLP thường đòi hỏi số lượng tham số lớn để đạt sai số nhỏ và khó đọc ra cấu trúc giải tích trực tiếp từ ma trận trọng số dày đặc (ngay cả trong các kiến trúc lớn như Transformer, khối MLP cũng chiếm phần lớn tham số phi nhúng~\cite{vaswani2017attention,liu2025kan}). Từ nhận xét đó, nhóm tác giả đề xuất \emph{Kolmogorov--Arnold Networks} (KAN) như một kiến trúc thay thế cho MLP bằng cách chuyển các hàm kích hoạt phi tuyến học được từ nút lên các cạnh nối.

Trong báo cáo này, cần phân biệt rõ hai vai trò dữ liệu: (i) ví dụ ứng dụng ở Mục~\ref{sec:motivating-example} được xây dựng để minh họa động lực thực tế của bài toán hồi quy phi tuyến nhiều biến; và (ii) các bài toán tổng hợp không nhiễu từ Mục~3.1 của bài báo~\cite{liu2025kan} được sử dụng trong thực nghiệm tái lập để đánh giá mối quan hệ giữa $\RMSE$ và số tham số.

\section{Ví dụ thực tế minh họa động lực}
\label{sec:motivating-example}
Để minh họa tầm quan trọng của việc xấp xỉ hàm nhiều biến với độ chính xác cao từ tập mẫu đo đạc vừa phải, xét một kịch bản ứng dụng thực tế giả định (đóng vai trò động lực nghiên cứu, không phải dữ liệu thực nghiệm của nhóm): \textbf{dự đoán công suất điện làm mát tức thời $P_{\mathrm{cool}}$ (kW) của một phòng máy chủ quy mô nhỏ} dựa trên hai đại lượng đo lường từ cảm biến môi trường gồm nhiệt độ không khí xung quanh $T \in [15, 40]\,^\circ\mathrm{C}$ và độ ẩm tương đối $H \in [20, 90]\,\%$. Tóm tắt bài toán được trình bày ở Bảng~\ref{tab:application}.

\begin{table}[htbp]\centering
\caption{Mô tả tình huống ứng dụng minh họa động lực bài toán xấp xỉ hàm.}\label{tab:application}
\begin{tabularx}{\linewidth}{@{}P{44mm}Y@{}}\toprule
Tình huống / người sử dụng & Kỹ sư vận hành phòng máy chủ quy mô nhỏ cần ước lượng tải điện làm mát khi điều kiện môi trường thay đổi.\\
Đầu vào và đơn vị & Véc-tơ $\mathbf{x} = (T, H)^\top \in \mathbb{R}^2$, gồm nhiệt độ $T$ ($^\circ\mathrm{C}$) và độ ẩm tương đối $H$ ($\%$).\\
Đầu ra cần dự đoán & Công suất điện làm mát $y = P_{\mathrm{cool}}$ (kW).\\
Quyết định được hỗ trợ & Lập lịch bật/tắt tổ máy lạnh dự phòng, phát hiện bất thường tiêu thụ điện và dự trù công suất nguồn trước khi nâng cấp cụm máy chủ.\\
Vì sao cần mô hình phi tuyến & Khi $T$ và $H$ tăng đồng thời, hiệu suất trao đổi nhiệt của dàn lạnh suy giảm phi tuyến và chịu tác động tương tác chéo giữa nhiệt độ điểm sương với tải nhiệt, khiến hồi quy tuyến tính bậc nhất không phản ánh đúng ngưỡng tăng công suất.\\\bottomrule
\end{tabularx}\end{table}

Trong kịch bản như trên, việc thu thập hàng triệu điểm đo có kiểm soát là tốn kém, trong khi kỹ sư vận hành lại cần mô hình gọn nhẹ, ít tham số và phản ánh được đóng góp riêng lẻ của từng biến đầu vào $T$ hay $H$. Đây chính là lớp bài toán mà nghiên cứu về KAN hướng tới~\cite{liu2025kan}.

\section{Ít nhất hai hướng tiếp cận trước KAN}
Trước khi KAN được đề xuất, bài toán xấp xỉ hàm liên tục nhiều biến chủ yếu dựa vào hai họ phương pháp tiêu biểu là mạng MLP và các phương pháp xấp xỉ bằng spline truyền thống~\cite{liu2025kan}.

\subsection{Multi-Layer Perceptron (MLP)}
\textbf{Ý tưởng cốt lõi:} MLP được truyền cảm hứng từ định lý xấp xỉ vạn năng~\cite{cybenko1989approximation,hornik1989multilayer,haykin1994neural}, cho phép một mạng truyền thẳng có ít nhất một lớp ẩn đủ rộng xấp xỉ tùy ý bất kỳ hàm liên tục nào trên miền compact. Về mặt hình thức, một mạng MLP gồm $L$ tầng biến đổi thực hiện xen kẽ phép biến đổi affine học được và hàm kích hoạt phi tuyến cố định tại từng nút~\cite{liu2025kan}:
\begin{equation}\label{eq:mlp-comp}
\mathrm{MLP}(\mathbf{x}) = \left(W_{L-1} \circ \sigma \circ W_{L-2} \circ \dots \circ \sigma \circ W_0\right)(\mathbf{x}),
\end{equation}
hay viết tường minh theo từng tầng $\ell \in \{0, \dots, L-1\}$:
\begin{equation}\label{eq:mlp-layer}
\mathbf{x}^{(\ell+1)} = \sigma\!\left(W^{(\ell)}\mathbf{x}^{(\ell)} + \mathbf{b}^{(\ell)}\right),
\end{equation}
trong đó $\mathbf{x}^{(0)} = \mathbf{x} \in \mathbb{R}^{n_0}$ là đầu vào, $W^{(\ell)} \in \mathbb{R}^{n_{\ell+1}\times n_\ell}$ là ma trận trọng số học được trên các cạnh, $\mathbf{b}^{(\ell)} \in \mathbb{R}^{n_{\ell+1}}$ là véc-tơ độ lệch (bias), còn $\sigma(\cdot)$ là hàm kích hoạt cố định áp dụng theo từng phần tử tại các nút (như $\mathrm{ReLU}$, $\tanh$ hoặc $\mathrm{SiLU}$; ở tầng đầu ra hồi quy thường bỏ $\sigma$). Điểm mạnh nổi bật của MLP là khả năng học đặc trưng hợp thành (feature learning) trong không gian nhiều chiều nhờ các phép chiếu tuyến tính $W^{(\ell)}$~\cite{liu2025kan}.

\textbf{Hạn chế và phạm vi kết luận:} Do tách rời biến đổi tuyến tính trên cạnh với hàm phi tuyến cố định tại nút, MLP không tối ưu hóa trực tiếp hình dạng của các hàm một biến. Trong phân tích của Liu và cộng sự~\cite{liu2025kan}, để xấp xỉ chính xác một hàm một biến trơn như $\sin(\pi x)$ hay $\exp(\cdot)$ chỉ bằng các hàm kích hoạt cố định, MLP buộc phải chồng nhiều nút ẩn và cần ngân sách tham số lớn. Theo phân tích quy luật tỷ lệ (scaling law) dựa trên xấp xỉ đa thức từng mảnh bậc $k$ trên đa tạp chiều nội tại $d$ của Sharma và Kaplan~\cite{sharma2020neural} (được trích dẫn đối chiếu trong~\cite{liu2025kan}), sai số kiểm tra của mạng dùng hàm kích hoạt cố định như ReLU ($k=1$) thường bị chặn tốc độ giảm quanh bậc $\alpha = (k+1)/d$, và thực nghiệm Mục~3.1 của~\cite{liu2025kan} cho thấy MLP có xu hướng bão hòa sớm khi tăng số tham số trên các hàm giải tích trơn. Ngoài ra, ma trận trọng số dày đặc $W^{(\ell)}$ khiến người dùng khó đọc ra công thức tường minh nếu không có công cụ giải thích hậu kỳ~\cite{liu2025kan}. Tuy nhiên, \emph{không được kết luận rằng MLP luôn kém hiệu quả hơn KAN trong mọi bài toán}: các nhận định về sự chậm bão hòa ở đây giới hạn trong lớp hàm giải tích trơn có cấu trúc hợp thành và thiết lập quy mô nhỏ của~\cite{liu2025kan}; trong các bài toán dữ liệu lớn và nhiều chiều như thị giác máy tính hoặc xử lý ngôn ngữ tự nhiên~\cite{vaswani2017attention}, MLP vẫn có lợi thế nổi bật về phép nhân ma trận song song trên phần cứng hiện đại và tốc độ huấn luyện thực tế.

\subsection{Xấp xỉ bằng spline truyền thống}
\textbf{Ý tưởng cốt lõi:} Trong giải tích số, một hàm một biến $s : [a, b] \to \mathbb{R}$ có thể được biểu diễn hiệu quả dưới dạng tổ hợp tuyến tính của các hàm cơ sở B-spline bậc $k$ định nghĩa trên lưới gồm $G$ khoảng chia~\cite{deboor1978practical,liu2025kan}:
\begin{equation}\label{eq:bspline-1d}
s(x) = \sum_{r=0}^{G+k-1} c_r B_r(x),
\end{equation}
trong đó $B_r(x)$ là hàm cơ sở B-spline thứ $r$ chỉ khác $0$ trên một khoảng giá đỡ cục bộ hẹp và $c_r \in \mathbb{R}$ là các hệ số điều khiển cần tìm. Nhờ tính chất cục bộ và cấu trúc đa thức từng mảnh trơn cấp $C^{k-1}$, spline một biến cho độ chính xác rất cao, dễ tinh chỉnh cục bộ mà không làm nhiễu vùng miền khác và dễ dàng tăng độ phân giải lưới~\cite{liu2025kan}.

\textbf{Hạn chế và phân biệt với spline trong KAN:} Khó khăn của spline cổ điển xuất hiện khi mở rộng trực tiếp lên không gian $d$ chiều bằng lưới tích ten-xơ (tensor-product spline). Khi đó, số lượng hàm cơ sở và hệ số tăng theo hàm mũ của số chiều, ở mức $\mathcal{O}\!\left((G+k)^d\right)$~\cite{liu2025kan}. Hệ quả là phương pháp lưới tích ten-xơ truyền thống chịu nặng \emph{lời nguyền số chiều} (curse of dimensionality -- COD) và không khai thác được cấu trúc hợp thành (compositional structure) tự nhiên của nhiều hàm thực tế (ví dụ $f(x_1,\dots,x_N)=\exp(\frac{1}{N}\sum_{i=1}^N \sin^2(x_i))$ trong Công thức~(1.1) của~\cite{liu2025kan}). Cần ghi nhận công bằng rằng tồn tại các biến thể spline nâng cao như sparse grids hoặc Generalized Additive Models (GAM) giúp giảm nhẹ số tham số dưới những giả định cấu trúc bổ sung. Điểm mấu chốt cần phân biệt là \textbf{KAN không sử dụng lưới spline đa chiều}: mỗi cạnh của KAN chỉ học một B-spline \emph{một biến} ($1\text{D}$), và tính đa biến được xây dựng bằng cách \emph{hợp thành qua các tầng nút cộng}, giữ tổng tham số tăng tuyến tính theo lưới ở mức $\mathcal{O}(LN^2(G+k))$ thay vì hàm mũ theo $d$.

\textbf{Cầu nối lịch sử tới Kolmogorov--Arnold:} Định lý biểu diễn Kolmogorov--Arnold (KAT)~\cite{kolmogorov1957representation} chứng minh mọi hàm liên tục nhiều biến trên miền bị chặn đều viết được dưới dạng hợp thành hữu hạn của phép cộng và các hàm một biến. Tuy nhiên, các nỗ lực trước đây khi áp dụng trực tiếp KAT vào mạng nơ-ron thường bị giới hạn ở kiến trúc 2 lớp cố định độ rộng $[d, 2d+1, 1]$ và vấp phải cảnh báo kinh điển của Girosi và Poggio~\cite{girosi1989representation} rằng các hàm một biến bên trong theo chứng minh gốc có thể không trơn (gây bất lợi cho học máy), dù về sau có các nghiên cứu chứng minh khả năng vượt COD khi giả định cấu trúc hợp thành~\cite{braun2009constructive,lai2021kolmogorov}. KAN~\cite{liu2025kan} kế thừa cả hai hướng trên: dùng cấu trúc nhiều lớp như MLP ở ``bên ngoài'' (external degrees of freedom) và tham số hóa các hàm một biến bằng B-spline ở ``bên trong'' (internal degrees of freedom). Bảng~\ref{tab:related} tổng hợp so sánh ba hướng tiếp cận.

\begin{table}[htbp]\centering\small
\caption{So sánh tổng quan giữa hai hướng tiếp cận trước đó và kiến trúc KAN theo~\cite{liu2025kan}.}\label{tab:related}
\begin{tabularx}{\linewidth}{@{}P{27mm}YYY@{}}\toprule
\textbf{Phương pháp} & \textbf{Ý tưởng cốt lõi} & \textbf{Hạn chế chính} & \textbf{Hệ quả / nguồn}\\\midrule
MLP & Phép biến đổi tuyến tính học được trên cạnh, kết hợp activation cố định tại nút~\eqref{eq:mlp-comp}. & Không tối ưu hóa trực tiếp hàm 1D; trọng số ma trận dày đặc khó giải thích trực tiếp. & Đường scaling tham số--sai số chậm hơn trong các bài toán giải tích nhỏ của~\cite{liu2025kan,sharma2020neural}.\\
Spline truyền thống (lưới ten-xơ) & Tổ hợp tuyến tính các B-spline giá đỡ cục bộ trên lưới $G$ khoảng~\eqref{eq:bspline-1d}. & Mở rộng trực tiếp lên $d$ chiều tốn $\mathcal{O}((G+k)^d)$ hệ số; không tự học cấu trúc hợp thành. & Bị lời nguyền số chiều (COD) khi $d$ tăng cao~\cite{deboor1978practical,liu2025kan}.\\
KAN & Tổng quát hóa KAT thành mạng sâu tùy ý; đặt hàm B-spline + SiLU học được lên cạnh, nút chỉ cộng. & Chi phí tính toán theo từng cạnh và tối ưu hóa với $G$ rất lớn phức tạp hơn phép nhân ma trận chuẩn. & Đạt đường scaling tốt hơn trên các hàm có cấu trúc trơn và hỗ trợ diễn giải~\cite{liu2025kan}.\\\bottomrule
\end{tabularx}\end{table}

\section{Đóng góp của bài báo và mục tiêu nhóm}
\label{sec:contributions-scope}
\subsection{Phân loại đóng góp của bài báo gốc}
Trong bài báo gốc~\cite{liu2025kan} (bản \codepath{arXiv:2404.19756v5}, được chấp nhận tại ICLR~2025), các đóng góp của tác giả có thể phân loại thành ba nhóm chính như sau:

\begin{enumerate}
\item \textbf{Đóng góp về lý thuyết (Mục~2.3, Định lý~2.1 của~\cite{liu2025kan}):} Thiết lập chặn xấp xỉ giải tích cho KAN nhiều lớp. Cụ thể, nếu hàm $f$ \emph{thừa nhận một biểu diễn KAN} trong đó mọi hàm cạnh đều khả vi liên tục cấp $k+1$, thì tồn tại xấp xỉ B-spline bậc $k$ trên lưới $G$ khoảng với sai số theo chuẩn $C^m$ ($0 \le m \le k$) bị chặn bởi $C\,G^{-(k+1-m)}$. \emph{Điểm mới so với các phương pháp trước:} Khác với định lý xấp xỉ vạn năng của MLP~\cite{hornik1989multilayer} vốn không cho tốc độ hội tụ miễn nhiễm với $d$, và khác với spline tích ten-xơ bị chi phí $\mathcal{O}((G+k)^d)$, chặn hội tụ $G^{-(k+1-m)}$ của KAN tách riêng độ tự do ngoài (cấu trúc hợp thành) và độ tự do trong (lưới $G$), không chứa $d$ ở số mũ (tuy hằng số $C$ vẫn phụ thuộc biểu diễn). Đồng thời cần lưu ý đây là \emph{chặn tồn tại xấp xỉ}, chưa tự động bảo đảm thuật toán tối ưu hội tụ tới nghiệm đó hay sai số test từ mẫu hữu hạn giảm đúng $G^{-4}$.
\item \textbf{Đóng góp về kỹ thuật (Mục~2.2, 2.4 và 2.5 của~\cite{liu2025kan}):} (i) Tổng quát hóa biểu diễn Kolmogorov--Arnold từ 2 lớp $[d, 2d+1, 1]$ lên kiến trúc sâu tùy ý $[n_0, n_1, \dots, n_L]$ huấn luyện bằng lan truyền ngược; (ii) kết hợp nhánh phần dư $\mathrm{SiLU}(x)$ với B-spline trên cạnh; (iii) kỹ thuật \emph{mở rộng lưới (grid extension)} giải bình phương tối thiểu khi tăng độ mịn lưới; và (iv) cơ chế thưa hóa ($L_1$/entropy), cắt tỉa nút (pruning), trực quan hóa và khớp hàm ký hiệu. \emph{Điểm mới so với các phương pháp trước:} Khắc phục rào cản không trơn và kiến trúc 2 lớp cố định của các mạng KAT cổ điển~\cite{girosi1989representation}, bổ sung khả năng nâng độ phân giải liên tiếp mà không phải huấn luyện lại từ đầu (điều MLP không làm trực tiếp trên kích hoạt nút).
\item \textbf{Đóng góp về thực nghiệm (Mục~3 và Mục~4 của~\cite{liu2025kan}):} Đối chứng KAN và MLP trên 5 hàm đồ chơi (Hình~3.1), 15 hàm đặc biệt (Mục~3.2), bộ phương trình Feynman (Mục~3.3), giải phương trình Poisson PDE (Mục~3.4), cùng hai bài toán khám phá khoa học (knot theory và chuyển pha định vị Anderson ở Mục~4). \emph{Điểm mới so với các phương pháp trước:} Cung cấp bằng chứng định lượng rằng trên các bài toán AI + Science quy mô nhỏ có cấu trúc hợp thành, KAN đạt đường scaling sai số--tham số dốc hơn MLP theo thực nghiệm và trực quan hóa được công thức giải tích thay vì làm ``hộp đen'' trọng số dày đặc.
\end{enumerate}

\subsection{Mục tiêu và phạm vi tái lập của nhóm}
Theo định hướng \textbf{tái lập cốt lõi} của đồ án môn học, nhóm lựa chọn tái lập sâu \textbf{một panel so sánh chính} từ Mục~3.1 / Hình~3.1 của bài báo~\cite{liu2025kan}:
\begin{equation}\label{eq:main-task}
f(x,y) = \exp\!\left(\sin(\pi x) + y^2\right),
\end{equation}
sử dụng kiến trúc KAN $[2, 1, 1]$ đối chứng trực tiếp với họ mô hình MLP trên cùng tập dữ liệu và phép đo $\RMSE$ theo số tham số học được.

Lý do lựa chọn phạm vi này là vì hàm~\eqref{eq:main-task} hội tụ đầy đủ cơ chế hợp thành phi tuyến 2 biến (biến đổi đơn biến $\sin(\pi x)$ và $y^2$, phép cộng tại nút giữa và hàm ngoài $\exp(\cdot)$), cho phép đánh giá rõ đường cong scaling KAN--MLP cũng như thực hiện \textbf{phân tích bổ sung bắt buộc về kích thước lưới $G$} (phân biệt lịch mở rộng lưới nối tiếp với mô hình khởi tạo lại ở từng $G$ cố định). Hai hàm bổ sung $J_0(20x)$ (KAN $[1,1]$) và $xy$ (KAN $[2,2,1]$) được giữ ở mức ưu tiên dự phòng nếu còn thời gian sau khi panel chính hoàn chỉnh. Các phần còn lại của bài báo gốc (15 hàm đặc biệt Mục~3.2, bộ dữ liệu Feynman Mục~3.3, giải PDE Mục~3.4 và hai thí nghiệm khám phá khoa học Mục~4) \textbf{không} được chạy lại mà chỉ được tóm tắt trong phần lý thuyết để bảo đảm độ sâu kiểm chứng trong quỹ thời gian 14 ngày.

\section{Dữ liệu tổng hợp và MLP đối chứng}
Phần này trình bày chi tiết quy cách sinh dữ liệu dùng chung và thiết lập các mô hình MLP đối chứng do TV1 phụ trách tại \codepath{Source/src/data.py} và \codepath{Source/src/mlp_baseline.py}. Bảng cấu hình khóa cuối cùng và kết quả thực nghiệm chính thức được tổng hợp đầy đủ tại Chương~\ref{ch:experiments}.

\subsection{Cơ chế dữ liệu và phân chia tập}
\label{sec:tv1-data-mechanism}
Khi đối chiếu nguồn bài báo~\cite{liu2025kan}, nhóm tách bạch rõ các thông số được xác định trực tiếp từ văn bản bài báo với các lựa chọn thiết kế của nhóm đối với những chi tiết không nêu tường minh trong PDF gốc:
\begin{itemize}
\item \textbf{Thông số xác định từ bài báo gốc:} Hàm mục tiêu chính $f(x,y)=\exp(\sin(\pi x)+y^2)$ và kiến trúc KAN $[2,1,1]$ tại Mục~3.1 (tr.~14); số lượng mẫu huấn luyện $n_{\mathrm{train}} = 1000$ cho bài toán này được ghi rõ ở Mục~2.4 (tr.~10: \emph{``our training samples are 1000''}).
\item \textbf{Thông số chưa xác định tường minh trong bài báo PDF và lựa chọn của nhóm:} Do bài báo không nêu cụ thể miền lấy mẫu của từng biến ở Mục~3.1, kích thước tập validation/test hay hạt giống ngẫu nhiên trong văn bản chính, nhóm cố định quy tắc chuẩn trong \codepath{Source/configs/dataset_exp_sin_y2.yaml}:
  \begin{enumerate}
  \item \textbf{Miền lấy mẫu và phương pháp lấy mẫu:} $\Omega = [-1, 1]^2$, lấy mẫu ngẫu nhiên phân phối đều liên tục $\mathcal{U}(-1, 1)$ độc lập cho từng tọa độ $(x, y)$.
  \item \textbf{Kích thước phân chia:} $n_{\mathrm{train}} = 1000$, $n_{\mathrm{val}} = 200$ (dùng riêng cho kiểm tra cấu hình/điểm dừng), và $n_{\mathrm{test}} = 1000$ (dùng chung cho đánh giá cuối).
  \item \textbf{Tính tách biệt và tái lập theo seed:} Từ hạt giống chung \codepath{seed = 42}, hàm \codepath{make_dataset()} sử dụng \codepath{np.random.SeedSequence(seed).spawn(3)} để sinh ba luồng RNG độc lập cho \codepath{train}, \codepath{val}, \codepath{test}. Toàn bộ tọa độ và nhãn mục tiêu được tính trước ở độ chính xác kép \codepath{float64} rồi mới chuyển sang \codepath{float32} nhằm giảm sai số làm tròn khi tính $\sin(\pi x)$ và $\exp(\cdot)$. Hàm tự động kiểm tra không có điểm trùng lặp giữa ba tập và khóa một mã định danh bất biến \codepath{split_id} (với \codepath{seed = 42}, mã kiểm chứng là \codepath{exp_sin_y2_s42_1000_200_1000_8150b7f5}).
  \item \textbf{Tiền xử lý, quy tắc chống rò rỉ và kiểm chứng giải tích:} Miền đầu vào $[-1, 1]^2$ đã đối xứng và có cùng thang đo trên hai trục, còn đầu ra $y \in [e^{-1}, e^{2}] \approx [0.37, 7.39]$ giữ nguyên đơn vị gốc để đối chiếu trực tiếp $\RMSE$ với Hình~3.1 của bài báo. Do đó nhóm cố định \codepath{normalize_x = False} và \codepath{normalize_y = False} trong \codepath{metadata} (TV2/TV3 không tự chuẩn hóa thêm ở pipeline sau). Nếu người dùng bật \codepath{normalize_x = True} cho bài toán mở rộng, thống kê trung bình/độ lệch chuẩn chỉ tính trên \codepath{x_train} rồi áp dụng cho \codepath{val} và \codepath{test}. Về vai trò tập: \codepath{train} dùng cho tối ưu gradient \codepath{LBFGS}, \codepath{val} dùng cho kiểm tra cấu hình/nhạy siêu tham số, còn \codepath{test} chỉ dùng đánh giá cuối. Hàm nạp split \codepath{DatasetBundle.load()} sử dụng \codepath{weights_only=True} để nạp an toàn dữ liệu nội bộ, và bộ kiểm thử \codepath{Source/tests/test_data.py} xác nhận sai số giải tích tại các điểm kiểm tra $f(0,0)=1$, $f(0,1)=e$, $f(1,0)=1$, $f(0.5,0)=e$, $f(-1,-1)=e$ nhỏ hơn $10^{-12}$.
  \end{enumerate}
\end{itemize}

\subsection{Baseline MLP, lựa chọn kích hoạt và phân tích thiết kế}
Để đối chứng theo cùng ngân sách tham số, nhóm xây dựng lớp \codepath{MLPBaseline} trong \codepath{Source/src/mlp_baseline.py} dựa trên PyTorch~\cite{paszke2019pytorch}, hỗ trợ cấu hình linh hoạt độ sâu, độ rộng, hàm kích hoạt (\codepath{silu}, \codepath{tanh}, \codepath{relu}, \codepath{gelu}), khởi tạo trọng số và kiểm soát \codepath{seed} độc lập qua \codepath{torch.random.fork_rng}.

\textbf{Lý giải lựa chọn kích hoạt $\mathrm{SiLU}$ cho baseline:} Cấu hình chính cố định $\sigma(z) = \mathrm{SiLU}(z) = z / (1 + e^{-z})$ ở các tầng ẩn, lớp đầu ra tuyến tính có \codepath{bias = True} và khởi tạo \codepath{xavier_uniform} vì hai lý do:
\begin{enumerate}
\item \emph{Tính trơn phù hợp với bài toán và bộ tối ưu \codepath{LBFGS}:} Hàm mục tiêu~\eqref{eq:main-task} khả vi vô hạn lần ($C^\infty$). Trong khi $\mathrm{ReLU}$ là hàm tuyến tính từng mảnh có điểm gãy không khả vi và đạo hàm bậc hai bằng $0$ hầu khắp nơi (làm chậm bước dò đường \codepath{strong_wolfe} của phương pháp tựa Newton \codepath{LBFGS}), $\mathrm{SiLU}$ thuộc lớp $C^\infty$ giúp gradient và xấp xỉ Hessian liên tục hơn.
\item \emph{Kiểm soát yếu tố nhiễu với KAN:} Nhánh cơ sở của KAN sử dụng $b(x)=\mathrm{SiLU}(x)$, nên dùng $\mathrm{SiLU}$ cho MLP giúp loại bớt sai khác do họ hàm kích hoạt khi so sánh kiến trúc.
\end{enumerate}
Để kiểm chứng nhận định trên, TV1 chạy kiểm tra độ nhạy phát triển (\codepath{Source/scripts/smoke_tv1.py}, 3~seed $\{42,43,44\}$, 200~vòng lặp \codepath{LBFGS} trên cấu hình trọng tâm 45~tham số \codepath{mlp_d2_w11}, lưu tại \codepath{fixtures/tv1_smoke_report.json}): với cùng ngân sách, $\mathrm{SiLU}$ ($\mathrm{val\_RMSE} = 0.1275 \pm 0.0320$, trung bình 1.72\,s) và $\mathrm{Tanh}$ ($0.1009 \pm 0.0247$, 1.79\,s) đều đạt sai số thấp hơn rõ rệt và hội tụ nhanh hơn $\mathrm{ReLU}$ ($0.2239 \pm 0.0185$, 2.61\,s). Ghi nhận trung thực rằng $\mathrm{Tanh}$ cho kết quả nhỉnh hơn $\mathrm{SiLU}$ ở thử nghiệm 200~bước này; đây là gợi ý quan trọng để nhóm kiểm tra thêm độ nhạy activation khi chạy đầy đủ 1800~bước trong Chương~\ref{ch:experiments}.

Với đầu vào 2 chiều và đầu ra 1 chiều có bias ở mọi lớp tuyến tính, số tham số học được (\codepath{requires_grad=True}) của MLP một lớp ẩn $[2, h, 1]$ là:
\begin{equation}\label{eq:mlp-params-1h}
P_{\mathrm{MLP}}([2, h, 1]) = (2h + h) + (h + 1) = 4h + 1,
\end{equation}
và của MLP hai lớp ẩn $[2, h, h, 1]$ là:
\begin{equation}\label{eq:mlp-params-2h}
P_{\mathrm{MLP}}([2, h, h, 1]) = (2h + h) + (h^2 + h) + (h + 1) = h^2 + 5h + 1.
\end{equation}
Sáu cấu hình MLP trong Bảng~\ref{tab:mlp-configs-locked} được \textbf{khóa trước hoàn toàn dựa trên công thức tham số} để phủ đều dải ngân sách dự kiến của KAN $[2,1,1]$ khi quét $G \in \{3, 5, 10, 20, 50\}$ trước khi quan sát bất kỳ giá trị $\RMSE$ nào trên tập kiểm tra. Lưu ý rằng \emph{sự tương đồng về số tham số chỉ bảo đảm so sánh trên cùng ngân sách tham số, không đồng nghĩa hai mô hình có cùng cấu trúc biểu diễn hay chi phí tính toán trên mỗi bước lặp}. Các chi tiết thiết kế kỹ thuật phần TV1 và tác động nếu thay đổi được tổng hợp ở Bảng~\ref{tab:tv1-code-decisions}.

\begin{table}[htbp]\centering\small
\caption{Sáu cấu hình MLP đối chứng được khóa trước theo số tham số trong \codepath{Source/configs/mlp_panel.yaml}.}\label{tab:mlp-configs-locked}
\begin{tabularx}{\linewidth}{@{}lccccY@{}}\toprule
\textbf{Config ID} & \textbf{Kiến trúc} & \textbf{Số lớp ẩn} & \textbf{Activation} & \textbf{Tham số} & \textbf{Vai trò đối chứng}\\\midrule
\codepath{mlp_d2_w6}  & $[2, 6, 1]$      & 1 & $\mathrm{SiLU}$ & 25  & Sát mức KAN $[2,1,1], G=3$\\
\codepath{mlp_d2_w11} & $[2, 11, 1]$     & 1 & $\mathrm{SiLU}$ & 45  & Sát mức KAN $G=10$ (cặp trọng tâm 3 seed)\\
\codepath{mlp_d2_w19} & $[2, 19, 1]$     & 1 & $\mathrm{SiLU}$ & 77  & Sát mức KAN $[2,1,1], G=20$\\
\codepath{mlp_d2_w41} & $[2, 41, 1]$     & 1 & $\mathrm{SiLU}$ & 165 & Sát mức KAN $[2,1,1], G=50$\\
\codepath{mlp_d3_w11} & $[2, 11, 11, 1]$ & 2 & $\mathrm{SiLU}$ & 177 & Kiểm tra ảnh hưởng độ sâu (depth~3) gần $G=50$\\
\codepath{mlp_d3_w20} & $[2, 20, 20, 1]$ & 2 & $\mathrm{SiLU}$ & 501 & Điểm mở rộng dung lượng ở depth~3\\\bottomrule
\end{tabularx}\end{table}

\begin{table}[htbp]\centering\small
\caption{Phân tích các lựa chọn thiết kế kỹ thuật trong mã nguồn TV1 (\codepath{data.py}, \codepath{mlp_baseline.py}).}\label{tab:tv1-code-decisions}
\begin{tabularx}{\linewidth}{@{}P{35mm}YY@{}}\toprule
\textbf{Thành phần mã nguồn} & \textbf{Cơ sở thiết kế} & \textbf{Tác động nếu sửa hoặc loại bỏ}\\\midrule
\codepath{SeedSequence.spawn(3)} & Tách 3 luồng RNG độc lập cho train/val/test. & Thay cỡ tập \codepath{n_train} sẽ làm dịch toàn bộ dãy mẫu \codepath{val/test}, mất tính đối chứng cùng tập test.\\
Tính nhãn bằng \codepath{float64} trước khi ép \codepath{float32} & Hạn chế sai số làm tròn khi tính $\sin(\pi x)$ và $\exp(\cdot)$. & Tính trực tiếp ở \codepath{float32} làm giảm độ chính xác nhãn chuẩn, đặc biệt nếu mở rộng xuống ngưỡng $\RMSE < 10^{-6}$.\\
\codepath{normalize_x/y = False} & Giữ nguyên miền $[-1,1]^2$ và thang đo $\RMSE$ gốc của bài báo. & Chuẩn hóa $y$ mà quên nghịch đảo khi đánh giá sẽ làm sai lệch hoàn toàn giá trị $\RMSE$ so với bài báo.\\
\codepath{fork_rng} khi khởi tạo MLP & Bảo đảm seed MLP tái lập nhưng không làm nhiễu RNG bên ngoài. & Bỏ cô lập RNG khiến trọng số khởi tạo của MLP phụ thuộc vào thứ tự gọi các hàm sinh dữ liệu trước đó.\\
Đếm tham số qua \codepath{requires_grad=True} & Chỉ đếm đúng các trọng số/bias thực sự được cập nhật. & Đếm cả buffer không học (như lưới knot của KAN) sẽ làm lệch trục số tham số khi vẽ đường scaling.\\\bottomrule
\end{tabularx}\end{table}

\reportfigure{figures/tv1/data_pipeline.pdf}{Sơ đồ quy trình dữ liệu, chia tập và huấn luyện đối chứng MLP--KAN.}{Quy trình sinh dữ liệu dùng chung, chia tập cố định và thiết lập đối chứng công bằng giữa MLP và KAN trên bài toán chính.}{fig:data-pipeline}
